% !TeX root = ../main.tex

\chapter{相关理论概述}



\section{卷积神经网络概述}

\subsection{卷积层}
\subsection{激活函数}
\subsection{池化层}
\subsection{全连接层}

\section{BYOL模型概述}

原版的BYOL模型处理对象为二维图像，目标是学习一个表示yθ，该表示随后可用于下游任务。BYOL使用两个神经网络进行学习：在线网络和目标网络。在线网络由一组权重 θ 定义，包含三个阶段：编码器（encoder）fθ，投影器（projector）gθ，预测器（predictor）qθ。目标网络具有与在线网络相同的架构，但使用另一组权重ξ。目标网络为训练在线网络提供回归目标（regression targets），其参数ξ是在线参数θ的指数移动平均值。更精确地说，给定一个目标衰减率（target decay rate）τ∈[0, 1]，在每次训练步骤后执行以下更新：ξ←τξ+(1-τ)θ。

输入与视图生成：给定一组图像D，一个从D中均匀采样的图像x，以及两个图像增强t和 t'。BYOL 通过分别应用图像增强t和 t'，从x生成两个增强视图：v ≜ t(x) 和 v' ≜ t'(x)。
前向传播与预测：从第一个增强视图v出发：在线网络输出一个表示yθ ≜ fθ(v)和一个投影 zθ ≜ gθ(y)。从第二个增强视图v'出发：目标网络输出一个表示 y'ξ ≜ fξ(v') 和一个目标投影 z'ξ ≜ gξ(y')。然后，在线网络输出目标投影 z'ξ 的一个预测 qθ(zθ)。预测器qθ仅应用于在线分支，这使得在线流程和目标流程之间的架构是不对称的。
损失函数：BYOL定义归一化后的预测与目标投影之间的均方误差为：
 
在每个训练步骤中执行随机优化步骤，仅针对参数θ最小化总损失.
参数更新：
BYOL 的动态更新过程总结如下：
 
训练结束时，仅保留编码器 fθ

\section{注意力机制概述}

\section{本章小结}

\section{论文章节安排}

\subsection{二级节标题}

\subsubsection{三级节标题}

本模板 \pkg{ustcthesis} 是中国科学技术大学本科生和研究生学位论文的 \LaTeX{}
模板, 按照《中国科学技术大学研究生学位论文撰写手册》（最近在修订中,以下简称《撰写手册》）和
《\href{https://www.teach.ustc.edu.cn/?attachment_id=13867}
{中国科学技术大学本科毕业论文（设计）格式}》的要求编写。

Lorem ipsum dolor sit amet, consectetur adipiscing elit, sed do eiusmod tempor
incididunt ut labore et dolore magna aliqua.
Ut enim ad minim veniam, quis nostrud exercitation ullamco laboris nisi ut
aliquip ex ea commodo consequat.
Duis aute irure dolor in reprehenderit in voluptate velit esse cillum dolore eu
fugiat nulla pariatur.
Excepteur sint occaecat cupidatat non proident, sunt in culpa qui officia
deserunt mollit anim id est laborum.


